\paragraph{\texorpdfstring{$\bar{\partial}$}{dbar}-算子的\trDeterminantLine。}%
\label{sub:determinant-lines-delbar}
我们需要回顾一些熟知的\trIndex{}理论结果，它们可以看作\eqref{eq:cut-decomposition}
的二维类似物。令 $S_k$（$k=1,2$）为两张带\trStriplikeEnd{}的\trRiemann\trSurface，
其上带有\trCauchyRiemann{}算子 $\bar{\partial}_{\nabla_k}$，按
\cref{sub:cauchy-riemann-operators} 的意义\trNondegenerate。设 $S_1$ 有一个出向端
$\zeta_1^{+}$，$S_2$ 有一个入向端 $\zeta_2^{-}$，并且两端上的\trLimitingDatum{}相
一致。于是可以把诸\trVectorBundle{}及其他数据\trGluing{}起来，得到\trSurface
$S=S_1\#_l S_2$（$l\gg 0$）上的一个\trNondegenerate{} \trCauchyRiemann{}算子
$\bar{\partial}_{\nabla}=\bar{\partial}_{\nabla_1}\#\bar{\partial}_{\nabla_2}$。
它在相差一个小\trPerturbation{}的意义下唯一；特别地，\trDeterminantLine{}
$\det(\bar{\partial}_{\nabla})$ 不依赖于构造中所做的选取。\trGluing{}定理说
\begin{equation}\label{eq:delbar-gluing}
\begin{gathered}
\mathrm{index}(\bar{\partial}_{\nabla})=\mathrm{index}(\bar{\partial}_{\nabla_1})+
\mathrm{index}(\bar{\partial}_{\nabla_2}),\\
\det(\bar{\partial}_{\nabla})\cong\det(\bar{\partial}_{\nabla_1})\otimes
\det(\bar{\partial}_{\nabla_2}).
\end{gathered}
\end{equation}
其中第二部分是一个\trCanonical{}\trIsomorphism，而且相对于双重\trGluing{}运算还是
结合的。证明\eqref{eq:delbar-gluing} 的标准方法是先假设两个参与算子都\trSurjective。
然后，对 $l\gg 0$，可以把诸 $\ker(\bar{\partial}_{\nabla_k})$ 中的元素拼补起来以逼近
$\bar{\partial}_{\nabla}=0$ 的解，这在经\trOrthogonal\trProjection{}之后给出一个%
\trIsomorphism
\begin{equation}\label{eq:delbar-kernel-iso}
\ker(\bar{\partial}_{\nabla})\cong\ker(\bar{\partial}_{\nabla_1})\oplus
\ker(\bar{\partial}_{\nabla_2});
\end{equation}
对极限 $l\to\infty$ 施以\trCompactness{}论证便表明 $\mathrm{coker}(\bar{\partial}_
{\nabla})=0$，所欲证的命题随之而来。\trSurjective{}性假设可以经一个有限维%
\trStabilization{}论证去掉。
偶尔我们也会用到另一种\trGluing{}步骤，它在 \trGromov--\trWitten{}理论中更为常见。
它有两个略微不同的版本，取决于\trGluing{}发生在\trInteriorPoint{}处还是边界点处，
但我们将同时讨论这两者。仍从两张带\trStriplikeEnd{}的\trSurface{} $S_k$ 出发，其上带有作用于%
\trVectorBundle{} $E_k$ 的算子 $\bar{\partial}_{\nabla_k}$，并有\trLagrange\trSubbundle
$F_k$。现在不是在诸端处取\trAsymptotic{}\trIsomorphism，而是设存在点
$z_k\in S_k$（要么同为\trInteriorPoint，要么同为边界点），以及\trFiber{}的一个等同关系
$E_{1,z_1}\cong E_{2,z_2}$，后者保持\trSymplecticStructure{}与\trComplexStructure{}（并且在边界的情形，还把诸%
\trLagrange{}子空间彼此映射）。作\trConnectedSum（若诸 $z_k$ 为\trInteriorPoint{}
则按通常意义；否则作\trBoundaryConnectedSum）便得到一张\trRiemann\trSurface{} $S$
以及其上的一个算子 $\bar{\partial}_{\nabla}$。\eqref{eq:delbar-kernel-iso} 的类似物
现在只对在诸点 $z_k$ 处取相同值的、由\trKernel{}元素组成的偶 $(X_1,X_2)$ 有定义，
而所得版本的\eqref{eq:delbar-gluing} 为：对\emph{边界点}，
\begin{equation}\label{eq:boundary-gluing}
\begin{gathered}
\mathrm{index}(\bar{\partial}_{\nabla})=\mathrm{index}(\bar{\partial}_{\nabla_1})+
\mathrm{index}(\bar{\partial}_{\nabla_2})-n,\\
\det(\bar{\partial}_{\nabla})\otimes\lambda^{\mathrm{top}}(F_{1,z_1})\cong
\det(\bar{\partial}_{\nabla_1})\otimes\det(\bar{\partial}_{\nabla_2})
\end{gathered}
\end{equation}
对\emph{\trInteriorPoint}，
\begin{equation}\label{eq:interior-gluing}
\begin{gathered}
\mathrm{index}(\bar{\partial}_{\nabla})=\mathrm{index}(\bar{\partial}_{\nabla_1})+
\mathrm{index}(\bar{\partial}_{\nabla_2})-2n,\\
\det(\bar{\partial}_{\nabla})\otimes\lambda^{\mathrm{top}}(E_{1,z_1})\cong
\det(\bar{\partial}_{\nabla_1})\otimes\det(\bar{\partial}_{\nabla_2}).
\end{gathered}
\end{equation}
可以利用 $E_{1,z_1}$ 由其\trSymplecticForm{}给出的\trOrientation，把 $\lambda^{\mathrm{top}}
(E_{1,z_1})$ 这一项从\eqref{eq:interior-gluing} 中去掉。还要注意，在\trInteriorPoint{}
情形，$S$ 与 $\bar{\partial}_{\nabla}$ 的构造依赖于一个额外的角度参数，它衡量诸%
\trTangentSpace{} $T(S_k)_{z_k}$ 的等同方式。然而，\eqref{eq:interior-gluing} 中的%
\trIsomorphism{}相对于该参数是连续的；因此，变动角度便给出算子 $\bar{\partial}_
{\nabla}$ 的一个 $S^1$-\trFamily，沿着它\trDeterminantLine{}\trLineBundle{}的单值
是平凡的。
\begin{remark}\label{thm:delbar-gluing-references}
诸如\eqref{eq:delbar-gluing} 之类的\trGluing{}公式出现在几乎每一个 \trFloerHomology
型理论的构造之中，例如见 [119]、[48]、[37, Section 3.3]。除此之外，关于一般%
\trDeterminantLine{}丛的切贴性质也有大量文献，其中很多基于 Atiyah--Patodi--Singer 型%
\trBoundaryCondition，而非（更简单却更受限的）管状端框架。导引性文献例如
[122, Lecture 2]；一般性定理见 [105]、[121]；\trRiemann\trSurface{}情形则见
[77, Appendix D]。
\end{remark}
